\subsection{13.01.2026 (Freistunde)}

\subsubsection{Ziel}

Ziel des Tages war es, die Skripte für die Ansteuerung des LCD-Displays
und die Erfassung der Sensordaten in das GitHub-Repository hochzuladen
sowie den Quellcode mithilfe von Klassen und Funktionen besser zu
strukturieren, um die Übersichtlichkeit und Wartbarkeit zu verbessern.

\subsubsection{Durchgeführte Arbeiten}

Zunächst wurden die Skripte \texttt{lcd.py} und
\texttt{dht\_11.py} in das GitHub-Repository hochgeladen.

Zusätzlich wurde das Skript \texttt{main.py} erstellt, welches
zukünftig die Hauptlogik der Taupunktlüftung steuern soll.

Das Skript \texttt{lcd.py} wurde anschließend in die Klasse
\texttt{LcdDisplay} umstrukturiert, um die Ansteuerung des
LCD-Displays zu kapseln und den Code übersichtlicher aufzubauen.

Weiterhin wurde der ursprüngliche Ablauf des Skripts als
Main-Methode innerhalb von \texttt{lcdDisplay.py} beibehalten,
damit die neue Struktur weiterhin getestet werden kann.

Darüber hinaus wurden die verschiedenen Aktionen des Skripts in
eigene Funktionen ausgelagert, um die Lesbarkeit und Wartbarkeit
des Codes zu verbessern. Diese Funktionen wurden anschließend
innerhalb der Main-Methode aufgerufen.

Ein praktischer Test der Änderungen war an diesem Tag jedoch nicht
möglich, da die JOY-PI-Hardware nicht zur Verfügung stand.

\subsubsection{Erkenntnisse}

Es wurde entschieden, die Skripte zukünftig nicht direkt vom
Raspberry Pi aus hochzuladen, sondern stattdessen einen lokalen
Computer für die Verwaltung des GitHub-Repositories zu verwenden.

Dadurch wird die Bearbeitung des Quellcodes erleichtert und das
Problem mit der GitHub-Authentifizierung auf dem Raspberry Pi
muss nicht weiter berücksichtigt werden.

Zusätzlich wurde festgestellt, dass eine Aufteilung des Programms
in Klassen und Funktionen die Übersichtlichkeit des Projekts
deutlich verbessert.

\subsubsection{Nächste Schritte}

In der nächsten Unterrichtseinheit sollen die überarbeiteten
Skripte auf dem JOY-PI getestet und weiterentwickelt werden.